% EvoClaw: A Self-Evolving Multi-Agent Framework
% with Genome-Driven Runtime Adaptation
%
% arXiv.cs.AI submission

\documentclass{article}
\usepackage[preprint]{neurips_2024}
\usepackage[utf8]{inputenc}
\usepackage[T1]{fontenc}
\usepackage{hyperref}
\usepackage{url}
\usepackage{booktabs}
\usepackage{amsfonts}
\usepackage{amsmath}
\usepackage{amssymb}

\title{EvoClaw: A Self-Evolving Multi-Agent Framework \\with Genome-Driven Runtime Adaptation}

\author{
  Alex Chen \\
  \texttt{alex.chen31337@gmail.com} \\
  \texttt{github.com/clawinfra/evoclaw}
}

\begin{document}

\maketitle

\begin{abstract}
Current multi-agent systems are fundamentally static: once deployed, an agent's architecture, capabilities, and behavior remain fixed until human engineers manually update them. This limitation prevents agents from adapting to dynamic environments, learning from experience, or evolving autonomously over time.

We present \textbf{EvoClaw}, a self-evolving multi-agent framework where agents modify their own behavior at runtime through a declarative genome specification. EvoClaw introduces three core innovations: (1) a \textit{genome.toml} file that declaratively specifies an agent's capabilities as pluggable skills; (2) a \textit{Genome Engine} that parses, validates, and hot-reloads skills without downtime; and (3) a \textit{multi-tier deployment model} supporting bare-metal edge devices, containerized servers, and cloud sandboxes. 

We demonstrate EvoClaw's practicality by deploying a fully functional agent on a Raspberry Pi 1 (ARMv6, 512MB RAM) running three concurrent skills—system monitoring, GPIO control, and cryptocurrency price tracking—with a binary size of 3.0MB and memory footprint of 3.2MB. Our Rust implementation achieves 90\%+ test coverage with 383 tests, and our Go orchestrator manages 100+ concurrent agents via MQTT pub/sub. 

EvoClaw enables a new class of autonomous agents that evolve without human intervention, opening research directions in self-modifying systems, edge AI, and agent economies. The complete system is open-source and reproducible.
\end{abstract}

\section{Introduction}

The rise of large language models (LLMs) has enabled autonomous agents that plan tasks, write code, and interact with systems \cite{auto_gpt,babyagi}. However, these agents suffer from a fundamental limitation: they are \textit{static}. Once deployed, an agent's architecture—its capabilities, behaviors, and response patterns—remains fixed until human engineers manually rewrite and redeploy it. This static architecture prevents agents from adapting to novel situations, learning from experience, or evolving autonomously over time.

We argue that truly autonomous agents should be able to modify themselves. Just as biological organisms evolve through natural selection, software agents should evolve through \textit{directed mutation}—guided by their own experience, environmental feedback, or human oversight. The challenge is designing a system that allows safe, verifiable self-modification without sacrificing reliability or performance.

\subsection{Contributions}

We present \textbf{EvoClaw}, a multi-agent framework where agents evolve through a declarative genome specification. Our contributions are:

\begin{itemize}
    \item \textbf{Genome-Driven Architecture}: We introduce \texttt{genome.toml}, a declarative configuration file that specifies an agent's capabilities as pluggable skills. Agents parse their genome at startup and can hot-reload when the genome changes.
    
    \item \textbf{Self-Evolving Rust Agents}: We implement a \textit{Genome Engine} in Rust that loads, validates, and executes skills dynamically. Skills are Rust traits with isolated execution, enabling zero-downtime updates.
    
    \item \textbf{Multi-Tier Deployment}: We design a deployment model supporting bare-metal edge devices (ARMv6, 512MB RAM), containerized servers (Podman), and cloud sandboxes (E2B). Agents synchronize across tiers via MQTT.
    
    \item \textbf{Production Deployment}: We demonstrate EvoClaw's practicality by deploying a fully-functional agent on a Raspberry Pi 1, running three concurrent skills (system monitoring, GPIO control, cryptocurrency price tracking) with a 3.0MB binary and 3.2MB memory footprint.
    
    \item \textbf{Open Source Release}: The complete system—including Go orchestrator (5,000 lines), Rust agent (8,000 lines), and documentation (30KB)—is released under an open-source license for full reproducibility.
\end{itemize}

\subsection{Paper Organization}

The remainder of this paper is organized as follows. Section \ref{sec:background} reviews related work in multi-agent systems and code synthesis. Section \ref{sec:architecture} describes EvoClaw's architecture. Section \ref{sec:genome} details the Genome Engine. Section \ref{sec:implementation} discusses implementation details. Section \ref{sec:evaluation} presents our evaluation on Raspberry Pi 1. Section \ref{sec:applications} describes applications. Section \ref{sec:limitations} discusses limitations and future work. Section \ref{sec:conclusion} concludes.

\section{Background}\label{sec:background}

\subsection{Multi-Agent Systems}

\textbf{Task Chaining Agents.} AutoGPT \cite{auto_gpt} and BabyAGI \cite{babyagi} decompose high-level goals into sequences of LLM-generated tasks. While powerful, these agents cannot modify their own architecture—new capabilities require human engineers to add code manually.

\textbf{Agent Frameworks.} LangChain \cite{langchain} provides abstractions for building LLM-powered agents (tools, memory, chains). However, agents are still static: changing behavior requires code changes and redeployment. CrewAI \cite{crewei} enables multi-agent collaboration but lacks runtime evolution.

\textbf{Edge AI.} TinyML \cite{tinyml} and edge computing frameworks optimize ML models for resource-constrained devices. However, these systems focus on inference, not agentic behavior or self-modification.

\subsection{Code Generation \& Synthesis}

\textbf{LLM Code Generation.} GPT-4 \cite{gpt4} and Claude \cite{claude} can write functional code, but they are not integrated into autonomous self-modifying systems. Program synthesis \cite{program_synthesis} focuses on generating code from specifications, not evolutionary adaptation.

\textbf{Self-Modifying Code.} Quines \cite{quine} and genetic programming \cite{gp} demonstrate self-replicating and evolving code, but these are theoretical constructs, not production multi-agent systems.

\subsection{Our Gap}

No existing system combines: (1) declarative agent specification, (2) hot-reloadable skills, (3) multi-tier deployment, and (4) production-ready edge deployment. EvoClaw bridges this gap.

\section{Architecture}\label{sec:architecture}

\subsection{System Overview}

Figure \ref{fig:architecture} shows EvoClaw's architecture. The system consists of three layers:

\begin{enumerate}
    \item \textbf{Go Orchestrator} (Director Mode): Manages agents, routes messages, and coordinates genome mutations. Provides a REST API and WebSocket interface for human interaction.
    
    \item \textbf{MQTT Broker}: Pub/sub message bus connecting orchestrator and agents. Supports topics like \texttt{agents/\{id\}/report} (telemetry) and \texttt{agents/\{id\}/commands} (genome updates).
    
    \item \textbf{Rust Edge Agents}: Lightweight binaries that execute skills defined in their \texttt{genome.toml}. Agents subscribe to MQTT topics and hot-reload when the genome changes.
\end{enumerate}

\subsection{Genome Structure}

Listing \ref{lst:genome} shows an example genome. The \texttt{[genome]} section defines agent identity and evolution parameters. The \texttt{[skills]} section declares enabled skills with configurations.

\begin{lstlisting}[language=toml, caption={Example \texttt{genome.toml}}, label={lst:genome}]
[genome]
id = "pi1-edge"
generation = 42
mutation_rate = 0.05

[skills.system_monitor]
enabled = true
interval_sec = 30

[skills.price_monitor]
enabled = true
symbols = ["BTC", "ETH", "SOL"]

[skills.gpio]
enabled = true
pins = [17, 27, 22]
\end{lstlisting}

\subsection{Multi-Tier Deployment}

EvoClaw supports three deployment tiers:

\begin{itemize}
    \item \textbf{Edge}: Bare-metal binary (ARMv6/ARMv7/x86\_64). 3.0MB with skills. Minimal resource usage (3.2MB RAM on Pi 1).
    
    \item \textbf{Server}: Podman containers with orchestrator. Scales to 100+ concurrent agents per orchestrator instance.
    
    \item \textbf{Cloud}: E2B sandboxes \cite{e2b} for isolated code execution. Useful for untrusted skills or testing.
\end{itemize}

Agents synchronize across tiers via MQTT, enabling hybrid deployments (e.g., edge agents report to cloud orchestrator).

\section{Genome Engine}\label{sec:genome}

\subsection{Skill Trait System}

Listing \ref{lst:trait} shows the \texttt{Skill} trait. Each skill implements \texttt{name()}, \texttt{execute()}, and \texttt{interval\_ms()}. Skills are isolated: errors in one skill do not crash the agent.

\begin{lstlisting}[language=rust, caption={Skill trait definition}, label={lst:trait}]
pub trait Skill: Send + Sync {
    fn name(&self) -> String;
    fn execute(&mut self, ctx: &Context) -> Result<Output>;
    fn interval_ms(&self) -> u64;
}
\end{lstlisting}

\subsection{Skill Registry}

The \texttt{SkillRegistry} maintains a map of skill names to instantiated trait objects. At startup, the registry:
\begin{enumerate}
    \item Parses \texttt{genome.toml}
    \item Instantiates enabled skills
    \item Validates configurations
    \item Spawns a task per skill that executes on a timer
\end{enumerate}

\subsection{Hot-Reload}

When the genome file changes, the agent:
\begin{enumerate}
    \item Detects the modification (via file watcher or MQTT signal)
    \item Re-parses \texttt{genome.toml}
    \item Computes the diff (added, removed, modified skills)
    \item Updates the registry (add/remove/reconfigure skills)
    \item Continues execution with zero downtime
\end{enumerate}

\section{Implementation}\label{sec:implementation}

\subsection{Technology Stack}

\textbf{Orchestrator (Go):} 5,000 lines across 11 packages. Uses MQTT (Eclipse Paho) for pub/sub, Gin for HTTP API, and libp2p for peer-to-peer discovery (planned).

\textbf{Edge Agents (Rust):} 8,000 lines with 383 tests (90\%+ coverage). Uses tokio for async runtime, serde for TOML parsing, and rumqtt for MQTT.

\textbf{Communication:} MQTT (Mosquitto broker) for lightweight pub/sub. Topics include \texttt{agents/\{id\}/report} (telemetry), \texttt{agents/\{id\}/commands} (genome updates), and \texttt{agents/\{id\}/chat} (human messages).

\subsection{Deployment Artifacts}

\begin{itemize}
    \item \texttt{evoclaw-orchestrator}: Go binary (7.2MB) for server deployment
    \item \texttt{evoclaw-agent}: Rust binary (3.0MB with skills) for edge devices
    \item \texttt{genome.toml}: Agent configuration
    \item \texttt{agent.toml}: Identity and credentials (MQTT, API keys)
\end{itemize}

\section{Evaluation}\label{sec:evaluation}

\subsection{Performance}

Table \ref{tab:performance} summarizes EvoClaw's performance.

\begin{table}[h]
\centering
\begin{tabular}{lr}
\toprule
\textbf{Metric} & \textbf{Value} \\
\midrule
Binary size (Rust agent) & 3.0 MB \\
RAM usage (Pi 1, 3 skills) & 3.2 MB \\
Skill load time & <100 ms \\
MQTT latency (pub/sub) & <50 ms \\
Test coverage (Rust) & 90\%+ \\
\bottomrule
\end{tabular}
\caption{EvoClaw performance metrics}
\label{tab:performance}
\end{table}

\subsection{Raspberry Pi 1 Deployment}

We deployed EvoClaw on a Raspberry Pi 1 (ARMv6, 512MB RAM) with three skills:
\begin{itemize}
    \item \textbf{SystemMonitor}: Reports CPU, RAM, disk, temperature every 30 seconds
    \item \textbf{GPIO}: Controls pins 17, 27, 22 (auto-detects Pi 1 GPIO offset of 512)
    \item \textbf{PriceMonitor}: Fetches BTC/ETH/SOL prices every 60 seconds
\end{itemize}

The agent ran for 24+ hours without restart, sending telemetry via MQTT. Total memory usage was 3.2MB RSS, well within the Pi's 512MB limit.

\subsection{Scalability}

The Go orchestrator manages 100+ concurrent agents via MQTT. Each agent maintains a lightweight TCP connection to the broker. The broker (Mosquitto) handles 10,000+ messages per second.

\section{Applications}\label{sec:applications}

\subsection{Edge Intelligence}

EvoClaw enables intelligent edge agents that adapt to local conditions. Example use cases:
\begin{itemize}
    \item IoT sensor networks (temperature, humidity, motion)
    \item Distributed computing (BOINC-style agent grids)
    \item Smart home automation (privacy-preserving, local processing)
\end{itemize}

\subsection{Cloud Automation}

EvoClaw agents run in E2B sandboxes for isolated code execution. Use cases:
\begin{itemize}
    \item Automated testing (CI/CD pipelines)
    \item Code review bots (analyze pull requests)
    \item Data processing (ETL workflows)
\end{itemize}

\subsection{Human-Agent Collaboration}

EvoClaw integrates with Telegram and Discord, enabling natural language interaction. Users can:
\begin{itemize}
    \item Query agent status via chat
    \item Request genome mutations (``add skill X'', ``remove Y'')
    \item Receive alerts and reports
\end{itemize}

\section{Limitations \& Future Work}\label{sec:limitations}

\subsection{Current Limitations}

\textbf{LLM Dependence.} Genome mutations are currently guided by LLMs (e.g., Claude, GPT-4). This requires external API calls and introduces latency/cost.

\textbf{No Formal Verification.} Evolved code is not formally verified, risking bugs or security issues. Safe mutation strategies are an open problem.

\textbf{Rust-Only Skills.} The current implementation only supports Rust skills. Cross-language support (e.g., Python, WASM) would broaden applicability.

\subsection{Future Directions}

\textbf{Formal Verification.} Integrate proof assistants (Coq, Lean) to verify evolved code preserves safety properties.

\textbf{WASM Skills.} Compile skills to WebAssembly for cross-language support and sandboxed execution.

\textbf{On-Chain Governance.} Integrate with ClawChain \cite{clawchain} for reputation-weighted genome mutation proposals.

\section{Conclusion}\label{sec:conclusion}

EvoClaw demonstrates that self-evolving agents are not just theoretically possible but practically deployable today. Our genome-based architecture enables autonomous adaptation across edge, server, and cloud environments—with production-ready deployment on resource-constrained hardware. 

By open-sourcing the complete system, we invite the community to extend EvoClaw's capabilities: new skill types, formal verification, and on-chain governance. The era of evolving agents has begun.

\section*{Acknowledgments}

Built with open-source tools: Rust, Go, MQTT, Substrate. Thanks to the Polkadot ecosystem for the Substrate framework.

\bibliographystyle{neurips_2024}
\begin{thebibliography}{9}

\bibitem{auto_gpt}
Toran Bruce Richards,
\newblock AutoGPT: An experimental open-source attempt to make GPT-4 fully autonomous,
\newblock 2023.

\bibitem{babyagi}
Yohei Nakajima,
\newblock babyagi: Managing tasks with AI,
\newblock 2023.

\bibitem{langchain}
Harrison Chase,
\newblock LangChain: Building applications with LLMs through composability,
\newblock 2022.

\bibitem{crewei}
Joao Moura,
\newblock CrewAI: Framework for orchestrating role-playing AI agents,
\newblock 2023.

\bibitem{tinyml}
\newblock TinyML: Machine learning on resource-constrained devices,
\newblock \url{https://www.tinyml.org}

\bibitem{e2b}
\newblock E2B: Type-safe TypeScript execution environments for AI agents,
\newblock \url{https://e2b.dev}

\bibitem{clawchain}
Alex Chen, Bowen Li,
\newblock ClawChain: An Agent-Native Layer 1 Blockchain,
\newblock 2026.

\end{thebibliography}

\end{document}
